\scanpage{41}\mspage{28}

Pour a), on applique $1^\circ$) \inserted{(ou simplement la \uline{déf.} de $C_I$)} dans le cas de l'inclusion d'ouverts $\varnothing\hookrightarrow I_{/i_{\unsure{0}}}$, \unsure{vu} qu'on a

\[
 u_{\varnothing,I_{/i_0}}\simeq u_{i_0}.
\]

Pour b), on applique $1^\circ$) à l'inclusion $\varnothing\longrightarrow I$, vu que

\[
 u_{\varnothing,I}\simeq u_I,
\]

et plus généralement

\[
 u_{\varnothing,J}\simeq u_J.
\]

\uline{Corollaire} 2.3. Supposons $u_\bullet\in C_I\cap W_{I^{\unsure{o|\circ}}}$. Alors \deleted{on} pour \unsure{tte} \deleted{\unsure{$i\in I$},} inclusion d'ouverts $K\subset J$ de $I$, on a $u_{\deleted{J}K,J}\deleted{\unsure{\text{est}}}\in TC$, et en particulier \deleted{(\unsure{3|2}.3.1)} (prenant $K=\varnothing$) on a $u_J\in TC$. Enfin pour tout $i\in I$, on a $u_i\in TC$.

\uline{Corollaire} 2.4. \deleted{\unsure{Supposons}} Soit $X_\bullet$ \unsure{de|dans} $M(I^o)_C$ un syst. inductif $C$-cofibrant, i.e. tel que $\varnothing_I\longrightarrow X_\bullet$ soit dans $C_I$, i.e. tel que \deleted{\unsure{les}} pour tout $i_0\in I$ et tout ouvert $K$ de $J=I_{/i_0}$, la flèche de transition $\tau_{K,J}:X_K=\underset{i\in K}{\underrightarrow{\lim}}\,X_i\longrightarrow X_J\simeq X_{i_0}$ soit dans $C$. Alors pour \uline{tout} couple d'ouverts $\deleted{J}K\subset\deleted{K}J$, la flèche \deleted{$\tau_{K,J}:$}

\[
 \tau_{K,J}:\underset{i\in K}{\unsure{\underrightarrow{\lim}|\underline{\lim}}}\,X_i
 \longrightarrow\underset{\unsure{i|i'}\in J}{\underrightarrow{\lim}}\,X_{\unsure{i'|i}}
\]

est dans $C$\textsuperscript{(*)}. \margin{et réciproquement, cela caractérise les objets $C$-cofibrants}
% slanted
En particulier (faisant $K=I_{/k}$, $J=I_{/j}$ pour $\unsure{h|k}\le j$) si $k\le j$ sont dans $I$, le morphisme de transition

\[
 \tau_{\unsure{k,j}}\quad X_k\unsure{\longrightarrow|-}X_j
\]

% archive sheet number: 40

